\documentclass[12pt]{article}
\usepackage[UTF8]{inputenc}
\usepackage{xeCJK}
\setCJKmainfont{Sarabun}
\usepackage{enumitem}
\usepackage{geometry}
\geometry{a4paper, margin=2.5cm}
\usepackage{setspace}
\onehalfspacing

\begin{document}

\begin{center}
{\Large\textbf{กัมมันตภาพรังสี}}\\[6pt]
{\normalsize วิชาฟิสิกส์ ม.ปลาย บทที่ 27}\\[4pt]
{\footnotesize ทั้งหมด 25 ข้อ}
\end{center}
\hrule
\bigskip

\section*{คำถาม in class}
\begin{enumerate}[label=\arabic*., itemsep=0.5\baselineskip, topsep=0.5\baselineskip, leftmargin=*]
\item นิวเคลียสของอะตอมประกอบด้วยอนุภาคใด?
\begin{enumerate}[label=\alph*)., itemsep=0.3\baselineskip, topsep=0.3\baselineskip]
\item ก\ 
\begin{minipage}[t]{0.9\linewidth}
\raggedright โปรตอนและอิเล็กตรอน
\end{minipage}
\item ข\ 
\begin{minipage}[t]{0.9\linewidth}
\raggedright โปรตอนและนิวตรอน
\end{minipage}
\item ค\ 
\begin{minipage}[t]{0.9\linewidth}
\raggedright นิวตรอนและอิเล็กตรออน
\end{minipage}
\item ง\ 
\begin{minipage}[t]{0.9\linewidth}
\raggedright โปรตอนเท่านั้น
\end{minipage}
\end{enumerate}
\vspace{-0.3\baselineskip}
\begin{small}
\textit{คำตอบ: ข}
\end{small}

\item ไอโซโทป (Isotope) คืออะตอมที่มีอะไรเหมือนกันแต่ต่างกัน?
\begin{enumerate}[label=\alph*)., itemsep=0.3\baselineskip, topsep=0.3\baselineskip]
\item ก\ 
\begin{minipage}[t]{0.9\linewidth}
\raggedright เลขอะตอมเท่ากัน แต่เลขมวลต่างกัน
\end{minipage}
\item ข\ 
\begin{minipage}[t]{0.9\linewidth}
\raggedright เลขมวลเท่ากัน แต่เลขอะตอมต่างกัน
\end{minipage}
\item ค\ 
\begin{minipage}[t]{0.9\linewidth}
\raggedright ทั้งเลขอะตอมและเลขมวลเท่ากัน
\end{minipage}
\item ง\ 
\begin{minipage}[t]{0.9\linewidth}
\raggedright จำนวนอิเล็กตรอนต่างกัน
\end{minipage}
\end{enumerate}
\vspace{-0.3\baselineskip}
\begin{small}
\textit{คำตอบ: ก}
\end{small}

\item การสลายตัวแบบแอลฟา (α) คืออะไร?
\begin{enumerate}[label=\alph*)., itemsep=0.3\baselineskip, topsep=0.3\baselineskip]
\item ก\ 
\begin{minipage}[t]{0.9\linewidth}
\raggedright การปล่อยอิเล็กตรอนจากนิวเคลียร์
\end{minipage}
\item ข\ 
\begin{minipage}[t]{0.9\linewidth}
\raggedright การปล่อยนิวเคลียสฮีเลียม (2p + 2n) ออกจากนิวเคลียร์
\end{minipage}
\item ค\ 
\begin{minipage}[t]{0.9\linewidth}
\raggedright การปล่อยโฟตอนพลังงานสูง
\end{minipage}
\item ง\ 
\begin{minipage}[t]{0.9\linewidth}
\raggedright การจับอิเล็กตรออน
\end{minipage}
\end{enumerate}
\vspace{-0.3\baselineskip}
\begin{small}
\textit{คำตอบ: ข}
\end{small}

\item การสลายตัวแบบเบต้า (β) แบ่งออกเป็นกี่ชนิด?
\begin{enumerate}[label=\alph*)., itemsep=0.3\baselineskip, topsep=0.3\baselineskip]
\item ก\ 
\begin{minipage}[t]{0.9\linewidth}
\raggedright 1 ชนิด
\end{minipage}
\item ข\ 
\begin{minipage}[t]{0.9\linewidth}
\raggedright 2 ชนิด
\end{minipage}
\item ค\ 
\begin{minipage}[t]{0.9\linewidth}
\raggedright 3 ชนิด
\end{minipage}
\item ง\ 
\begin{minipage}[t]{0.9\linewidth}
\raggedright 4 ชนิด
\end{minipage}
\end{enumerate}
\vspace{-0.3\baselineskip}
\begin{small}
\textit{คำตอบ: ข}
\end{small}

\item การสลายตัวแบบแกมมา (γ) คืออะไร?
\begin{enumerate}[label=\alph*)., itemsep=0.3\baselineskip, topsep=0.3\baselineskip]
\item ก\ 
\begin{minipage}[t]{0.9\linewidth}
\raggedright การปล่อยอิเล็กตรออน
\end{minipage}
\item ข\ 
\begin{minipage}[t]{0.9\linewidth}
\raggedright การปล่อยโพซิตรอน
\end{minipage}
\item ค\ 
\begin{minipage}[t]{0.9\linewidth}
\raggedright การปล่อยโฟตอนพลังงานสูงจากนิวเคลียร์
\end{minipage}
\item ง\ 
\begin{minipage}[t]{0.9\linewidth}
\raggedright การปล่อยนิวตรอน
\end{minipage}
\end{enumerate}
\vspace{-0.3\baselineskip}
\begin{small}
\textit{คำตอบ: ค}
\end{small}

\item ครึ่งชีวิต (half-life) ของสารกัมมันตรังสีคืออะไร?
\begin{enumerate}[label=\alph*)., itemsep=0.3\baselineskip, topsep=0.3\baselineskip]
\item ก\ 
\begin{minipage}[t]{0.9\linewidth}
\raggedright เวลาที่สารสลายตัวจนหมด
\end{minipage}
\item ข\ 
\begin{minipage}[t]{0.9\linewidth}
\raggedright เวลาที่จำนวนนิวเคลียสลดลงเหลือครึ่งหนึ่งของจำนวนเริ่มต้น
\end{minipage}
\item ค\ 
\begin{minipage}[t]{0.9\linewidth}
\raggedright เวลาที่สารกัมมันตรังสีเปลี่ยนสถานะ
\end{minipage}
\item ง\ 
\begin{minipage}[t]{0.9\linewidth}
\raggedright เวลาที่อะตอมเปลี่ยนธาตุ
\end{minipage}
\end{enumerate}
\vspace{-0.3\baselineskip}
\begin{small}
\textit{คำตอบ: ข}
\end{small}

\item กฎการสลายตัวของสารกัมมันตรังสี N = N₀e\textasciicircum{}(-λt) โดย λ คืออะไร?
\begin{enumerate}[label=\alph*)., itemsep=0.3\baselineskip, topsep=0.3\baselineskip]
\item ก\ 
\begin{minipage}[t]{0.9\linewidth}
\raggedright ครึ่งชีวิต
\end{minipage}
\item ข\ 
\begin{minipage}[t]{0.9\linewidth}
\raggedright ค่าคงที่การสลายตัว (decay constant)
\end{minipage}
\item ค\ 
\begin{minipage}[t]{0.9\linewidth}
\raggedright กัมมันตภาพรังสี
\end{minipage}
\item ง\ 
\begin{minipage}[t]{0.9\linewidth}
\raggedright จำนวนนิวเคลียสเริ่มต้น
\end{minipage}
\end{enumerate}
\vspace{-0.3\baselineskip}
\begin{small}
\textit{คำตอบ: ข}
\end{small}

\item การหาอายุด้วยคาร์บอน-14 ใช้หลักการอะไร?
\begin{enumerate}[label=\alph*)., itemsep=0.3\baselineskip, topsep=0.3\baselineskip]
\item ก\ 
\begin{minipage}[t]{0.9\linewidth}
\raggedright เปลี่ยน C-14 เป็น C-12
\end{minipage}
\item ข\ 
\begin{minipage}[t]{0.9\linewidth}
\raggedright เปรียบเทียบอัตราส่วน C-14/C-12 ในซากดึกดำบรรพ์กับอัตราส่วนในอากาศ
\end{minipage}
\item ค\ 
\begin{minipage}[t]{0.9\linewidth}
\raggedright นับจำนวนอะตอมทั้งหมด
\end{minipage}
\item ง\ 
\begin{minipage}[t]{0.9\linewidth}
\raggedright วัดน้ำหนักของซากดึกดำบรรพ์
\end{minipage}
\end{enumerate}
\vspace{-0.3\baselineskip}
\begin{small}
\textit{คำตอบ: ข}
\end{small}

\item พลังงานยึดเหนี่ยวนิวเคลียร์ (nuclear binding energy) คืออะไร?
\begin{enumerate}[label=\alph*)., itemsep=0.3\baselineskip, topsep=0.3\baselineskip]
\item ก\ 
\begin{minipage}[t]{0.9\linewidth}
\raggedright พลังงานที่ใช้ในปฏิกิริยานิวเคลียร์
\end{minipage}
\item ข\ 
\begin{minipage}[t]{0.9\linewidth}
\raggedright พลังงานที่ใช้ในการแยกนิวเคลียร์ออกเป็นโปรตอนและนิวตรอน
\end{minipage}
\item ค\ 
\begin{minipage}[t]{0.9\linewidth}
\raggedright พลังงานที่ยึดนิวตรอนไว้ในอิเล็กตรออน
\end{minipage}
\item ง\ 
\begin{minipage}[t]{0.9\linewidth}
\raggedright พลังงานจลน์ของนิวเคลียร์
\end{minipage}
\end{enumerate}
\vspace{-0.3\baselineskip}
\begin{small}
\textit{คำตอบ: ข}
\end{small}

\item อนุภาค β⁻ ที่ปล่อยออกมาจากการสลายตัวแบบเบต้าลบ คืออะไร?
\begin{enumerate}[label=\alph*)., itemsep=0.3\baselineskip, topsep=0.3\baselineskip]
\item ก\ 
\begin{minipage}[t]{0.9\linewidth}
\raggedright นิวเคลียร์ฮีเลียม
\end{minipage}
\item ข\ 
\begin{minipage}[t]{0.9\linewidth}
\raggedright อิเล็กตรออน
\end{minipage}
\item ค\ 
\begin{minipage}[t]{0.9\linewidth}
\raggedright โพซิตรอน
\end{minipage}
\item ง\ 
\begin{minipage}[t]{0.9\linewidth}
\raggedright โฟตอน
\end{minipage}
\end{enumerate}
\vspace{-0.3\baselineskip}
\begin{small}
\textit{คำตอบ: ข}
\end{small}

\item สารกัมมันตรังสีมีครึ่งชีวิต 6 ชั่วโมง ถ้าเริ่มต้นมี 80 g หลังผ่านไป 18 ชั่วโมง จะเหลือกี่กรัม?
\begin{enumerate}[label=\alph*)., itemsep=0.3\baselineskip, topsep=0.3\baselineskip]
\item ก\ 
\begin{minipage}[t]{0.9\linewidth}
\raggedright 40 g
\end{minipage}
\item ข\ 
\begin{minipage}[t]{0.9\linewidth}
\raggedright 20 g
\end{minipage}
\item ค\ 
\begin{minipage}[t]{0.9\linewidth}
\raggedright 10 g
\end{minipage}
\item ง\ 
\begin{minipage}[t]{0.9\linewidth}
\raggedright 5 g
\end{minipage}
\end{enumerate}
\vspace{-0.3\baselineskip}
\begin{small}
\textit{คำตอบ: ค}
\end{small}

\item ค่าคงที่การสลายตัว λ ของไอโซโทปที่มีครึ่งชีวิต 10 วัน มีค่าเท่าใด?
\begin{enumerate}[label=\alph*)., itemsep=0.3\baselineskip, topsep=0.3\baselineskip]
\item ก\ 
\begin{minipage}[t]{0.9\linewidth}
\raggedright 0.0693 วัน⁻¹
\end{minipage}
\item ข\ 
\begin{minipage}[t]{0.9\linewidth}
\raggedright 0.0693 ชม.⁻¹
\end{minipage}
\item ค\ 
\begin{minipage}[t]{0.9\linewidth}
\raggedright 6.93 วัน⁻¹
\end{minipage}
\item ง\ 
\begin{minipage}[t]{0.9\linewidth}
\raggedright 10 วัน⁻¹
\end{minipage}
\end{enumerate}
\vspace{-0.3\baselineskip}
\begin{small}
\textit{คำตอบ: ก}
\end{small}

\item เมื่อนิวเคลียส Y-90 (Z=39, A=90) สลายตัวโดยปล่อย β⁻ ผลที่ได้คือนิวเคลียสใหม่ที่มี Z และ A เท่าใด?
\begin{enumerate}[label=\alph*)., itemsep=0.3\baselineskip, topsep=0.3\baselineskip]
\item ก\ 
\begin{minipage}[t]{0.9\linewidth}
\raggedright Z=40, A=90
\end{minipage}
\item ข\ 
\begin{minipage}[t]{0.9\linewidth}
\raggedright Z=38, A=90
\end{minipage}
\item ค\ 
\begin{minipage}[t]{0.9\linewidth}
\raggedright Z=40, A=91
\end{minipage}
\item ง\ 
\begin{minipage}[t]{0.9\linewidth}
\raggedright Z=39, A=89
\end{minipage}
\end{enumerate}
\vspace{-0.3\baselineskip}
\begin{small}
\textit{คำตอบ: ก}
\end{small}

\item ครึ่งชีวิตของคาร์บอน-14 เท่ากับประมาณกี่ปี?
\begin{enumerate}[label=\alph*)., itemsep=0.3\baselineskip, topsep=0.3\baselineskip]
\item ก\ 
\begin{minipage}[t]{0.9\linewidth}
\raggedright 1,400 ปี
\end{minipage}
\item ข\ 
\begin{minipage}[t]{0.9\linewidth}
\raggedright 5,730 ปี
\end{minipage}
\item ค\ 
\begin{minipage}[t]{0.9\linewidth}
\raggedright 14,000 ปี
\end{minipage}
\item ง\ 
\begin{minipage}[t]{0.9\linewidth}
\raggedright 50,000 ปี
\end{minipage}
\end{enumerate}
\vspace{-0.3\baselineskip}
\begin{small}
\textit{คำตอบ: ข}
\end{small}

\item ไอโซโทปของยูเรเนียมที่ใช้เป็นเชื้อเพลิงในโรงไฟฟ้านิวเคลียร์คืออะไร?
\begin{enumerate}[label=\alph*)., itemsep=0.3\baselineskip, topsep=0.3\baselineskip]
\item ก\ 
\begin{minipage}[t]{0.9\linewidth}
\raggedright U-235
\end{minipage}
\item ข\ 
\begin{minipage}[t]{0.9\linewidth}
\raggedright U-238
\end{minipage}
\item ค\ 
\begin{minipage}[t]{0.9\linewidth}
\raggedright U-234
\end{minipage}
\item ง\ 
\begin{minipage}[t]{0.9\linewidth}
\raggedright U-239
\end{minipage}
\end{enumerate}
\vspace{-0.3\baselineskip}
\begin{small}
\textit{คำตอบ: ก}
\end{small}

\item กัมมันตภาพรังสี A (activity) ของสารกัมมันตรังสี มีความสัมพันธ์กับจำนวนนิวเคลียส N อย่างไร?
\begin{enumerate}[label=\alph*)., itemsep=0.3\baselineskip, topsep=0.3\baselineskip]
\item ก\ 
\begin{minipage}[t]{0.9\linewidth}
\raggedright A = N/λ
\end{minipage}
\item ข\ 
\begin{minipage}[t]{0.9\linewidth}
\raggedright A = λN
\end{minipage}
\item ค\ 
\begin{minipage}[t]{0.9\linewidth}
\raggedright A = N₀e\textasciicircum{}(-λt)
\end{minipage}
\item ง\ 
\begin{minipage}[t]{0.9\linewidth}
\raggedright A = N₀ - N
\end{minipage}
\end{enumerate}
\vspace{-0.3\baselineskip}
\begin{small}
\textit{คำตอบ: ข}
\end{small}

\item รังสีใดมีอำนาจทะลุทะลวงสูงที่สุด?
\begin{enumerate}[label=\alph*)., itemsep=0.3\baselineskip, topsep=0.3\baselineskip]
\item ก\ 
\begin{minipage}[t]{0.9\linewidth}
\raggedright แอลฟา (α)
\end{minipage}
\item ข\ 
\begin{minipage}[t]{0.9\linewidth}
\raggedright เบต้า (β)
\end{minipage}
\item ค\ 
\begin{minipage}[t]{0.9\linewidth}
\raggedright แกมมา (γ)
\end{minipage}
\item ง\ 
\begin{minipage}[t]{0.9\linewidth}
\raggedright ทุกตัวมีอำนาจทะลุทะลวงเท่ากัน
\end{minipage}
\end{enumerate}
\vspace{-0.3\baselineskip}
\begin{small}
\textit{คำตอบ: ค}
\end{small}

\item หลังจาก 3 ครึ่งชีวิต สารกัมมันตรังสีจะเหลือกี่เปอร์เซ็นต์ของจำนวนเริ่มต้น?
\begin{enumerate}[label=\alph*)., itemsep=0.3\baselineskip, topsep=0.3\baselineskip]
\item ก\ 
\begin{minipage}[t]{0.9\linewidth}
\raggedright 50\%
\end{minipage}
\item ข\ 
\begin{minipage}[t]{0.9\linewidth}
\raggedright 25\%
\end{minipage}
\item ค\ 
\begin{minipage}[t]{0.9\linewidth}
\raggedright 12.5\%
\end{minipage}
\item ง\ 
\begin{minipage}[t]{0.9\linewidth}
\raggedright 6.25\%
\end{minipage}
\end{enumerate}
\vspace{-0.3\baselineskip}
\begin{small}
\textit{คำตอบ: ค}
\end{small}

\item นิวเคลียร์ธาตุเหล็ก-56 (²⁶Fe⁵⁶) มีจำนวนโปรตอน นิวตรอน และอิเล็กตรออนเท่าใด?
\begin{enumerate}[label=\alph*)., itemsep=0.3\baselineskip, topsep=0.3\baselineskip]
\item ก\ 
\begin{minipage}[t]{0.9\linewidth}
\raggedright 26p, 26n, 26e
\end{minipage}
\item ข\ 
\begin{minipage}[t]{0.9\linewidth}
\raggedright 26p, 30n, 26e
\end{minipage}
\item ค\ 
\begin{minipage}[t]{0.9\linewidth}
\raggedright 26p, 26n, 30e
\end{minipage}
\item ง\ 
\begin{minipage}[t]{0.9\linewidth}
\raggedright 56p, 26n, 56e
\end{minipage}
\end{enumerate}
\vspace{-0.3\baselineskip}
\begin{small}
\textit{คำตอบ: ข}
\end{small}

\item นิวเคลียส Rn-222 (Z=86) สลายตัวโดยปล่อย α จะได้นิวเคลียสใหม่ที่มี Z และ A เท่าใด?
\begin{enumerate}[label=\alph*)., itemsep=0.3\baselineskip, topsep=0.3\baselineskip]
\item ก\ 
\begin{minipage}[t]{0.9\linewidth}
\raggedright Z=84, A=220
\end{minipage}
\item ข\ 
\begin{minipage}[t]{0.9\linewidth}
\raggedright Z=84, A=222
\end{minipage}
\item ค\ 
\begin{minipage}[t]{0.9\linewidth}
\raggedright Z=88, A=222
\end{minipage}
\item ง\ 
\begin{minipage}[t]{0.9\linewidth}
\raggedright Z=88, A=226
\end{minipage}
\end{enumerate}
\vspace{-0.3\baselineskip}
\begin{small}
\textit{คำตอบ: ข}
\end{small}

\item สารกัมมันตรังสี 10 mg มีครึ่งชีวิต 5 วัน หลัง 20 วันจะเหลือกี่มิลลิกรัม?
\begin{enumerate}[label=\alph*)., itemsep=0.3\baselineskip, topsep=0.3\baselineskip]
\item ก\ 
\begin{minipage}[t]{0.9\linewidth}
\raggedright 5.0 mg
\end{minipage}
\item ข\ 
\begin{minipage}[t]{0.9\linewidth}
\raggedright 2.5 mg
\end{minipage}
\item ค\ 
\begin{minipage}[t]{0.9\linewidth}
\raggedright 1.25 mg
\end{minipage}
\item ง\ 
\begin{minipage}[t]{0.9\linewidth}
\raggedright 0.625 mg
\end{minipage}
\end{enumerate}
\vspace{-0.3\baselineskip}
\begin{small}
\textit{คำตอบ: ง}
\end{small}

\item ซากดึกดำบรรพ์มี C-14 เหลือ 25\% ของปกติ ซากดึกดำบรรพ์นี้มีอายุประมาณเท่าใด? (T½ C-14 = 5730 ปี)
\begin{enumerate}[label=\alph*)., itemsep=0.3\baselineskip, topsep=0.3\baselineskip]
\item ก\ 
\begin{minipage}[t]{0.9\linewidth}
\raggedright 5730 × 2 = 11460 ปี
\end{minipage}
\item ข\ 
\begin{minipage}[t]{0.9\linewidth}
\raggedright 5730 × 3 = 17190 ปี
\end{minipage}
\item ค\ 
\begin{minipage}[t]{0.9\linewidth}
\raggedright 5730 × 4 = 22920 ปี
\end{minipage}
\item ง\ 
\begin{minipage}[t]{0.9\linewidth}
\raggedright 5730 × 1.5 = 8595 ปี
\end{minipage}
\end{enumerate}
\vspace{-0.3\baselineskip}
\begin{small}
\textit{คำตอบ: ค}
\end{small}

\item เมื่อโปรตอนในนิวเคลียร์เปลี่ยนเป็นนิวตรอน จะปล่อยอะไรออกมา?
\begin{enumerate}[label=\alph*)., itemsep=0.3\baselineskip, topsep=0.3\baselineskip]
\item ก\ 
\begin{minipage}[t]{0.9\linewidth}
\raggedright อิเล็กตรออน
\end{minipage}
\item ข\ 
\begin{minipage}[t]{0.9\linewidth}
\raggedright โพซิตรอน + นิวตริโน
\end{minipage}
\item ค\ 
\begin{minipage}[t]{0.9\linewidth}
\raggedright แอลฟา
\end{minipage}
\item ง\ 
\begin{minipage}[t]{0.9\linewidth}
\raggedright แกมมา
\end{minipage}
\end{enumerate}
\vspace{-0.3\baselineskip}
\begin{small}
\textit{คำตอบ: ข}
\end{small}

\item สารกัมมันตรังสีมีค่าคงที่การสลายตัว λ = 0.01 ชม.⁻¹ ครึ่งชีวิตเท่ากับเท่าใด?
\begin{enumerate}[label=\alph*)., itemsep=0.3\baselineskip, topsep=0.3\baselineskip]
\item ก\ 
\begin{minipage}[t]{0.9\linewidth}
\raggedright 0.693 ชม.
\end{minipage}
\item ข\ 
\begin{minipage}[t]{0.9\linewidth}
\raggedright 6.93 ชม.
\end{minipage}
\item ค\ 
\begin{minipage}[t]{0.9\linewidth}
\raggedright 69.3 ชม.
\end{minipage}
\item ง\ 
\begin{minipage}[t]{0.9\linewidth}
\raggedright 100 ชม.
\end{minipage}
\end{enumerate}
\vspace{-0.3\baselineskip}
\begin{small}
\textit{คำตอบ: ค}
\end{small}

\item นิวเคลียสกัมมันตรังสีมีจำนวนเริ่มต้น 10¹² นิวเคลียส และมี λ = 0.1 วินาที⁻¹ กัมมันตภาพรังสีเริ่มต้นเท่ากับเท่าใด?
\begin{enumerate}[label=\alph*)., itemsep=0.3\baselineskip, topsep=0.3\baselineskip]
\item ก\ 
\begin{minipage}[t]{0.9\linewidth}
\raggedright 10¹¹ Bq
\end{minipage}
\item ข\ 
\begin{minipage}[t]{0.9\linewidth}
\raggedright 10¹² Bq
\end{minipage}
\item ค\ 
\begin{minipage}[t]{0.9\linewidth}
\raggedright 10¹³ Bq
\end{minipage}
\item ง\ 
\begin{minipage}[t]{0.9\linewidth}
\raggedright 10¹⁰ Bq
\end{minipage}
\end{enumerate}
\vspace{-0.3\baselineskip}
\begin{small}
\textit{คำตอบ: ก}
\end{small}

\end{enumerate}
\end{document}